\chapter{How to Make \$2 Million in Commercial Real Estate}

This lecture is not a blackboard lecture in the usual sense. Its mathematics is hidden inside prices, rent streams, bid rules, deal clocks, and commission structures. We follow the speaker in the order he gives the material: first the large transactions, then the path into the business, then the daily discipline, then the claim that commercial brokerage is built less on isolated transactions than on repeated relationships. Where the interview only suggests a formula, we keep the reconstruction modest and say so.

\section{Opening Setup}

Let us begin where the lecture begins: with the biggest deal, then with the correction that reorganizes the opening. The first answer is a \$45 million retail property in New York City. Then, after a pause, the speaker remembers a larger and more interesting deal. That self-interruption matters. The lecture does not proceed from theory to example; it proceeds from example to better example, and only then to method.

Before we return to those deals, the speaker resets the discussion and explains how he entered the business at all. In 1999 he joins a New York City firm with no experience, after basketball ambitions give way to restaurant work and then to the discovery that real estate is not merely about one's own home but about the ownership and transfer of buildings as an asset class. The first part of the lecture therefore asks us to hold two facts at once: the beginner knows almost nothing, and the market he enters is already a world of very large numbers.

\begin{equation}
2009 - 1999 = 10 \text{ years}
\end{equation}

That simple chronology is worth recording because the speaker stays ten years at the first firm before leaving in 2009 to start his own company. The lecture does not present success as immediate improvisation. It presents it as apprenticeship followed by duration.

\section{Two Deals and the Logic of Selection}

The speaker returns to the first headline transaction: a retail condo in New York City, occupied by HSBC Bank, with annual rental income of about \$2.5 million and a final deal size of about \$45 million. He also says that the conversation at one stage began around \$55 million, and that he regarded that level as overpriced. No formal valuation formula is spoken in the interview, but we can write down the smallest reconstruction that captures the intuition.

\begin{remark}
The equations in this section are transcript-derived summaries. They are not board equations from the lecture, and they should be read as compact ways of expressing the speaker's pricing judgment.
\end{remark}

\begin{align}
P_{\text{deal}} &\approx \$45\,\text{M}, &
R &\approx \$2.5\,\text{M/yr}, &
y &= \frac{R}{P}.
\end{align}

\paragraph{A worked example.}
If we use gross rent yield only as a first-pass comparison, then the same rent stream supports different pricing stories at \$45 million and \$55 million.

\begin{enumerate}
\item Start with the quoted rent \(R \approx \$2.5\,\text{M/yr}\).
\item Evaluate the ratio \(y = R/P\) at the final deal price \(P=45\).
\item Compare that value with the same ratio at the earlier discussed level \(P=55\).
\item Use the contrast only as a cautious explanation of the remark that \$55 million felt too high.
\end{enumerate}

\begin{align}
y_{45} &\approx \frac{2.5}{45} \approx 0.0556 \approx 5.6\%, \\
y_{55} &\approx \frac{2.5}{55} \approx 0.0455 \approx 4.5\%.
\end{align}

The point is not that the interview offers a full underwriting model. The point is simpler: price is not just a headline. It is a relation between income and acquisition cost, and a few percentage points in that relation can separate an attractive deal from one that feels too rich.

Now the lecture sharpens. The speaker remembers the larger transaction: a Bronx portfolio of apartment buildings handled through the U.S. Marshals. Here the speaker does something more severe than give us another large number. He states a rule. It is a closed bid process, and the winner is simply the highest offer, even if that offer is higher by only one cent. The buyer's ability to close does not enter the selection rule.

\begin{definition}
In the narrow sense relevant here, a closed bid process is one in which each bidder submits a number \(b_i\), and the winner is determined solely by the largest submitted bid.
\end{definition}

\begin{equation}
i^\* = \arg\max_i b_i
\end{equation}

The lecture pushes this rule to its extreme form:

\begin{equation}
b_i = b_j + \$0.01 \quad \Longrightarrow \quad i \text{ wins over } j.
\end{equation}

That is the mechanism. It strips away all comforting stories about superior judgment after the bid is in. If the Marshals' rule is exactly as described, then the decisive object is the submitted number, not the narrative around the number.

\begin{figure}[h]
\centering
\begin{tikzpicture}[x=1cm,y=1cm]
\node[draw, minimum width=2.2cm, minimum height=0.8cm] (b1) at (0,2.0) {$b_1$};
\node[draw, minimum width=2.2cm, minimum height=0.8cm] (b2) at (0,0.8) {$b_2$};
\node[draw, minimum width=2.2cm, minimum height=0.8cm] (b3) at (0,-0.4) {$b_3$};
\node[draw, minimum width=3.8cm, minimum height=1.2cm] (m) at (6,0.8) {U.S. Marshals\\closed bid rule};
\draw[->] (b1) -- (m);
\draw[->] (b2) -- (m);
\draw[->] (b3) -- (m);
\node at (6,-1.0) {$i^\*=\arg\max_i b_i$};
\node at (6,-1.8) {highest offer wins};
\end{tikzpicture}
\end{figure}

The speaker then supplies the outcome: the bid that won was about \$73,500,550, though he marks the number as approximate. We should preserve that approximation rather than pretending to a false exactness. What matters most for the chapter is the pedagogical function of the anecdote: the lecture has moved from valuation to market design.

\section{The Discipline of the Brokerage Regime}

Once the two large deals are on the table, the lecture pivots away from glamour and toward routine. The question becomes: what is the day-to-day life of someone running a brokerage? The answer is an operating regime. The day begins around 5:30, arrival at the office is by 6:15 at the latest, departure is around 7:30 or 8:00 at night, and the week is structured by recurring meetings and reviews.

This part of the lecture is easy to underrate, but it performs a necessary correction. We have just heard about deals in the tens of millions. Now we are told that the underlying machine is repetitive, managerial, and scheduled. Weekly meetings ask what went well, what did not go well, and how business can be built more effectively in the next cycle. The lecture thus replaces the mythology of the single heroic transaction with the more sober picture of a process repeated over decades.

A useful way to read this section is to notice that the speaker has now supplied three different time scales. There is the daily clock of work, the monthly clock of pipeline development, and the multi-year clock on which reputations and earnings compound. The later quantitative claims make little sense unless we keep all three clocks in view at once.

\section{Relationships, No Salary, and the Carrying-Cost Problem}

The first explicit tip for someone trying to break into the industry is also the governing thesis of the lecture: commercial brokers are not transactional; they are relationship-centric. One does not ideally close a single deal and disappear. One returns to the same clients, again and again, over long stretches of time.

The speaker gives a striking heuristic. If one can land about five active clients, one is effectively set.

\begin{equation}
N_{\text{active}} \approx 5
\end{equation}

We should not treat this as a theorem. It is a rule of thumb offered by a manager to a young broker. But it clarifies the structure of the business. The problem is not to accumulate hundreds of shallow contacts. The problem is to establish a small enough set of relationships that each can recur and deepen.

The lecture now introduces the harsh economic condition that makes the whole business intelligible: there is no salary. The speaker's response to that condition is not fear but appetite. No salary means no ceiling. But the optimistic half of that statement depends on surviving the pessimistic half. If commissions are delayed, then a beginner must somehow carry himself through the barren interval before the first closings arrive.

The speaker tells us how he did it. He worked long brokerage hours during the day, then took a restaurant job at night earning roughly \$10 to \$12 an hour, and also sold cars on weekends. The lecture refuses the easy moral here. It does not universalize the path or present exhaustion as virtue for its own sake. Instead it identifies a structural fact: a commission business with long deal cycles imposes carrying costs on the entrant, and those costs must be borne somehow if the entrant is to stay in the game long enough for relationships to mature.

\section{First Deal, First Commission, First Capital}

At this point the interview turns toward framed plaques from earlier transactions. We do not need the image of the wall to understand the logic. What matters is the first recorded deal and what it did in the speaker's financial life.

The first transaction is dated August 22, 2001, in Binghamton, New York. It took almost two years to close. The commission check was \$54,000.

\begin{equation}
C_1 = \$54{,}000
\end{equation}

That figure is not introduced as a boast. It functions as a bridge between labor and capital. The lecture's claim is that commissions from these early apartment-building transactions later became the basis for the speaker's first personal deal. In other words, brokerage fees are not only income. They can also become acquisition capital.

This is one of the most important transitions in the chapter. The lecture began with brokerage as intermediary work. Here it reveals a second layer. Intermediation can become ownership, but only after delay. First comes the pipeline, then the first commission, then repeated commissions, and only then the possibility of buying on one's own account.

\section{Networking as an Access Problem}

The lecture now returns to relationships, but at a more precise level. The question is no longer whether relationships matter. That has already been decided. The question is what someone should do if networking does not come naturally.

\subsection{Question \& Answer}

\paragraph{Question.}
What if one struggles with networking and does not naturally build lasting relationships?

\paragraph{Answer.}
The lecture answers by reframing networking as attention rather than performance. One begins from one's own strong points, studies the individual one is trying to reach, and looks for a real point of contact rather than a generic sales script.

The example is the speaker's attempt to reach Julian Studley, who would later become a mentor. Phone calls produced no response. Letters produced no response. Progress begins only when the speaker discovers that Studley is a poker player. He then opens the conversation there, not by pretending expertise, but by saying in effect: I hear you are a poker expert; I would like to learn.

We can summarize the method as a sequence.

\begin{enumerate}
\item Identify the person whose attention matters.
\item Gather information about what genuinely matters to that person.
\item Use that information to open a sincere conversation.
\item Convert a single meeting into repeated contact.
\item Let repeated contact mature into business.
\end{enumerate}

The lecture's proof of concept is immediate and concrete. After joining poker nights and learning the game, the speaker is hired to sell about \$30 million in real estate. The point is not that poker is magical. The point is that access is often gained by understanding the person's world before asking for a transaction.

\section{Licensing, Transition, and the Time-to-Income Curve}

The interview now raises a different practical obstacle. What is the process of going to work in commercial real estate, and should one begin in residential and then move over later? The lecture answers both questions by focusing on time.

First, the licensing point. The speaker says the licensing portion is the same. One goes to the same real-estate school whether the intended practice is residential or commercial. The real distinction is not at the level of the license itself. It is at the level of business structure, deal duration, and the ability to tolerate delay.

\subsection{Question \& Answer}

\paragraph{Question.}
What is the actual path into commercial brokerage, and is residential real estate a natural staging ground?

\paragraph{Answer.}
The lecture's answer has three pieces. The licensing school is the same. The commercial side is much slower. Because it is slower, starting in residential is not obviously a good preparation if commercial brokerage is the real aim.

The timing claims can be written compactly.

\begin{align}
T_{\text{first}}^{\text{com}} &\in [6,12]\text{ months}, \\
T_{\text{close}}^{\text{res}} &\in [30,60]\text{ days}, \\
T_{\text{close}}^{\text{com}} &\approx 120\text{ days}.
\end{align}

These numbers explain much of the lecture at once. If residential deals can close in as little as a month or two, or even within a week when a friend or cousin is ready to transact, then residential cash flow can begin quickly. Commercial brokerage is different. The first commercial close can take six to twelve months, and even after listing a property the path to closing is still long.

That is why the speaker does not recommend a casual transition from residential into commercial. The habits of the two businesses differ because the clocks differ. Residential can reward an agent early. Commercial forces the broker to build a pipeline that can survive long gaps.

The lecture then pushes from the time-to-income curve toward long-run earnings claims. By year five, if one is doing the business correctly, building a database, building relationships, presenting proposals, and becoming an expert, the claim is that one can be in seven-figure territory. After a decade, the claim is roughly \$10 million or more.

\begin{align}
Y(5\text{ years}) &\in \text{seven figures}, \\
Y(10\text{ years}) &\gtrsim \$10\,\text{M}.
\end{align}

We should read these as the speaker's claims about the upside of the field, not as laws of nature. The lecture itself supplies the necessary check: when asked about his own best year, the speaker says it was about \$2 million.

\begin{equation}
Y_{\text{best}} \approx \$2\,\text{M}
\end{equation}

That answer is important. It brings the lecture back down from horizon talk to lived experience. It also opens the final observation about technology. The speaker says that today's brokers have access to information that earlier brokers did not: search tools, owner data, and fast discovery. Information, in his phrase, is gold. What earlier generations had to obtain by knocking on doors, later entrants can often surface from a database.

\section{Summary}

The lecture's mathematics is modest in form but sharp in consequence. A property's income must be compared to its price. A closed bid process selects the highest bid, not the best story. A commission business without salary imposes carrying costs on the entrant. Commercial deal cycles are much longer than residential ones, and that difference governs both early hardship and later upside.

But the lecture does not want us to stop at formulas. Its central thesis is that commercial brokerage is a long-horizon game of repeated relationships. The first task is entry, the second is endurance, the third is repetition, and only then do large commissions and personal ownership become realistic. That is the real sequence of the chapter, and it is also the real sequence of the lecture.